\subsection{Diagonalizable Operators}

\begin{exercise}[5d1]
    Suppose $V$ is a finite-dimensional complex vector space and $T \in \lin{V}$.
    \begin{subexercises}
        \item Prove that if $T^4 = I$, then $T$ is diagonalizable.
        \item Prove that if $T^4 = T$, then $T$ is diagonalizable.
        \item Give an example of an operator $T \in \lin{\cbb^2}$ such that $T^4 = T^2$ and $T$ is not diagonalizable.
    \end{subexercises}
    \remark{The roots $\omega$ in part (b) are the roots obtained in \exref{1c6}.}
\end{exercise}

\begin{sole}
    \item If $T^4 = I$, then the minimal polynomial of $T$ divides $z^4 - 1 = (z - 1)(z + 1)(z - i)(z + i)$, which has no repeated roots. Hence the minimal polynomial of $T$ has no repeated roots, so $T$ is diagonalizable.
    \item If $T^4 = T$, then the minimal polynomial of $T$ divides $z^4 - z = z(z - 1)(z - \omega)(z - \omega^2)$. This polynomial has no repeated roots, so the minimal polynomial of $T$ has no repeated roots, and hence $T$ is diagonalizable.
    \item Let $T \in \lin{\cbb^2}$ be defined by
    \[
        T(z_1, z_2) = (z_2, 0).
    \]
    Then $T^2(z_1, z_2) = (0, 0)$, so $T^4 = T^2$. However, $T$ is not diagonalizable because the only eigenvalue of $T$ is $0$, and $\dim E(0, T) = 1 < 2 = \dim \cbb^2$.
\end{sole}

\begin{exercise}[5d2]
    Suppose $T \in \lin{V}$ has a diagonal matrix $A$ with respect to some basis of $V$.
    Prove that if $\lambda \in \fbb$, then $\lambda$ appears on the diagonal of $A$ precisely $\dim E(\lambda, T)$ times.
\end{exercise}

\begin{sol}
    Suppose $A$ is a diagonal matrix of $T$ with respect to some basis of $V$.
    Let $\lambda \in \fbb$.
    Let $m$ denote the number of times $\lambda$ appears on the diagonal of $A$.
    Then $A - \lambda I$ is a diagonal matrix of $T - \lambda I$ with respect to the same basis, with exactly $m$ zeros on the diagonal.
    A diagonal matrix with $m$ zeros on the diagonal has null space of dimension $m$, so $\dim \nul(T - \lambda I) = m$.
    Since $E(\lambda, T) = \nul(T - \lambda I)$, we have $\dim E(\lambda, T) = m$.
\end{sol}

\begin{exercise}[5d3]
    Suppose $V$ is finite-dimensional and $T \in \lin{V}$.
    Prove that if the operator $T$ is diagonalizable, then $V = \nul T \op \range T$.
\end{exercise}

\begin{sol}
    Suppose $T$ is diagonalizable.
    Let $\lambda_1, \ldots, \lambda_m$ denote the distinct eigenvalues of $T$.
    Then
    \[
        V = E(\lambda_1, T) \op \cdots \op E(\lambda_m, T).
    \]
    If 0 is not an eigenvalue of $T$, then $\nul T = \set{0}$ and $\range T = V$, so $V = \nul T \op \range T$.
    If 0 is an eigenvalue of $T$, then we may assume $\lambda_1 = 0$.
    Let
    \[
        U = E(\lambda_2, T) \op \cdots \op E(\lambda_m, T).
    \]
    Since $\nul T = E(0, T)$, we have $V = \nul T \op U$.
    It remains to show that $U = \range T$.

    Let $u \in U$.
    Then $u = u_2 + \cdots + u_m$ for some $u_k \in E(\lambda_k, T)$.
    Since $\lambda_k \neq 0$ for $k = 2, \ldots, m$, we have $u_k = \lambda_k\inv T(u_k)$ for each $k = 2, \ldots, m$.
    Hence,
    \[
        u = u_2 + \cdots + u_m = \lambda_2\inv T(u_2) + \cdots + \lambda_m\inv T(u_m) = T(\lambda_2\inv u_2 + \cdots + \lambda_m\inv u_m) \in \range T.
    \]
    Thus, $U \subseteq \range T$.

    Now let $v \in \range T$.
    Note that eigenspaces are invariant as follows: If $u \in E(\lambda_k, T)$, then $T(u) = \lambda_k u \in E(\lambda_k, T)$.
    Since $U$ is a direct sum of eigenspaces, it is invariant under $T$.
    Then $v = T(w)$ for some $w \in V$.
    Since $V = \nul T \op U$, we may write $w = w_1 + w_2$ for some $w_1 \in \nul T$ and $w_2 \in U$.
    Then
    \[
        v = T(w) = T(w_1 + w_2) = T(w_1) + T(w_2) = 0 + T(w_2) = T(w_2) \in U.
    \]
    Thus, $\range T \subseteq U$.
    Since $U \subseteq \range T$ and $\range T \subseteq U$, we have $U = \range T$.
    Therefore, $V = \nul T \op U = \nul T \op \range T$.
\end{sol}

\begin{exercise}[5d4]
    Suppose $V$ is finite-dimensional and $T \in \lin{V}$.
    Prove that the following are equivalent.
    \begin{subexercises}
        \item $V = \nul T \op \range T$.
        \item $V = \nul T + \range T$.
        \item $\nul T \cap \range T = \set{0}$.
    \end{subexercises}
\end{exercise}

\begin{exercise}[5d5]
    Suppose $V$ is a finite-dimensional complex vector space and $T \in \lin{V}$.
    Prove that $T$ is diagonalizable if and only if
    \[
        V = \nul(T - \lambda I) \op \range(T - \lambda I)
    \]
    for every $\lambda \in \cbb$.
\end{exercise}

\begin{exercise}[5d6]
    Suppose $T \in \lin{\fbb^5}$ and $\dim E(8, T) = 4$.
    Prove that $T - 2I$ or $T - 6I$ is invertible.
\end{exercise}

\begin{exercise}[5d7]
    Suppose $T \in \lin{V}$ is invertible.
    Prove that
    \[
        E(\lambda, T) = E\left(\tfrac{1}{\lambda}, T\inv\right)
    \]
    for every $\lambda \in \fbb$ with $\lambda \neq 0$.
\end{exercise}

\begin{exercise}[5d8]
    Suppose $V$ is finite-dimensional and $T \in \lin{V}$.
    Let $\lambda_1, \ldots, \lambda_m$ denote the distinct nonzero eigenvalues of $T$.
    Prove that
    \[
        \dim E(\lambda_1, T) + \cdots + \dim E(\lambda_m, T) \leq \dim \range T.
    \]
\end{exercise}

\begin{exercise}[5d9]
    Suppose $R, T \in \lin{\fbb^3}$ each have $2, 6, 7$ as eigenvalues.
    Prove that there exists an invertible operator $S \in \lin{\fbb^3}$ such that $R = S\inv T S$.
\end{exercise}

\begin{exercise}[5d10]
    Find $R, T \in \lin{\fbb^4}$ such that $R$ and $T$ each have $2, 6, 7$ as eigenvalues, $R$ and $T$ have no other eigenvalues, and there does not exist an invertible operator $S \in \lin{\fbb^4}$ such that $R = S\inv T S$.
\end{exercise}

\begin{exercise}[5d11]
    Find $T \in \lin{\cbb^3}$ such that $6$ and $7$ are eigenvalues of $T$ and such that $T$ does not have a diagonal matrix with respect to any basis of $\cbb^3$.
\end{exercise}

\begin{exercise}[5d12]
    Suppose $T \in \lin{\cbb^3}$ is such that $6$ and $7$ are eigenvalues of $T$.
    Furthermore, suppose $T$ does not have a diagonal matrix with respect to any basis of $\cbb^3$.
    Prove that there exists $(z_1, z_2, z_3) \in \cbb^3$ such that
    \[
        T(z_1, z_2, z_3) = (6 + 8z_1, 7 + 8z_2, 13 + 8z_3).
    \]
\end{exercise}

\begin{exercise}[5d13]
    Suppose $A$ is a diagonal matrix with distinct entries on the diagonal and $B$ is a matrix of the same size as $A$.
    Show that $AB = BA$ if and only if $B$ is a diagonal matrix.
\end{exercise}

\begin{exercise}[5d14]
    \begin{subexercises}
        \item Give an example of a finite-dimensional complex vector space and an operator $T$ on that vector space such that $T^2$ is diagonalizable but $T$ is not diagonalizable.
        \item Suppose $\fbb = \cbb$, $k$ is a positive integer, and $T \in \lin{V}$ is invertible.
        Prove that $T$ is diagonalizable if and only if $T^k$ is diagonalizable.
    \end{subexercises}
\end{exercise}

\begin{exercise}[5d15]
    Suppose $V$ is a finite-dimensional complex vector space, $T \in \lin{V}$, and $p$ is the minimal polynomial of $T$.
    Prove that the following are equivalent.
    \begin{subexercises}
        \item $T$ is diagonalizable.
        \item There does not exist $\lambda \in \cbb$ such that $p$ is a polynomial multiple of $(z - \lambda)^2$.
        \item $p$ and its derivative $p'$ have no zeros in common.
        \item The greatest common divisor of $p$ and $p'$ is the constant polynomial $1$.
    \end{subexercises}
    \textit{The \emph{\textbf{greatest common divisor}} of $p$ and $p'$ is the monic polynomial $q$ of largest degree such that $p$ and $p'$ are both polynomial multiples of $q$.
    The Euclidean algorithm for polynomials (look it up) can quickly determine the greatest common divisor of two polynomials, without requiring any information about the zeros of the polynomials.
    Thus the equivalence of (a) and (d) above shows that we can determine whether $T$ is diagonalizable without knowing anything about the zeros of $p$.}
\end{exercise}

\begin{exercise}[5d16]
    Suppose that $T \in \lin{V}$ is diagonalizable.
    Let $\lambda_1, \ldots, \lambda_m$ denote the distinct eigenvalues of $T$.
    Prove that a subspace $U$ of $V$ is invariant under $T$ if and only if there exist subspaces $U_1, \ldots, U_m$ of $V$ such that $U_k \subseteq E(\lambda_k, T)$ for each $k$ and $U = U_1 \op \cdots \op U_m$.
\end{exercise}

\begin{exercise}[5d17]
    Suppose $V$ is finite-dimensional.
    Prove that $\lin{V}$ has a basis consisting of diagonalizable operators.
\end{exercise}

\begin{exercise}[5d18]
    Suppose that $T \in \lin{V}$ is diagonalizable and $U$ is a subspace of $V$ that is invariant under $T$.
    Prove that the quotient operator $T/U$ is a diagonalizable operator on $V/U$.
\end{exercise}

\begin{exercise}[5d19]
    Prove or give a counterexample: If $T \in \lin{V}$ and there exists a subspace $U$ of $V$ that is invariant under $T$ such that $T|_U$ and $T/U$ are both diagonalizable, then $T$ is diagonalizable.
    \textit{See Exercise 13 in Section 5C for an analogous statement about upper-triangular matrices.}
\end{exercise}

\begin{exercise}[5d20]
    Suppose $V$ is finite-dimensional and $T \in \lin{V}$.
    Prove that $T$ is diagonalizable if and only if the dual operator $T'$ is diagonalizable.
\end{exercise}

\begin{exercise}[5d21]
    The \emph{\textbf{Fibonacci sequence}} $F_0, F_1, F_2, \ldots$ is defined by
    \[
        F_0 = 0, \quad F_1 = 1, \quad \text{and } F_n = F_{n-2} + F_{n-1} \text{ for } n \geq 2.
    \]
    Define $T \in \lin{\rbb^2}$ by $T(x, y) = (y, x + y)$.
    \begin{subexercises}
        \item Show that $T^n(0, 1) = (F_n, F_{n+1})$ for each nonnegative integer $n$.
        \item Find the eigenvalues of $T$.
        \item Find a basis of $\rbb^2$ consisting of eigenvectors of $T$.
        \item Use the solution to (c) to compute $T^n(0, 1)$.
        Conclude that
        \[
            F_n = \frac{1}{\sqrt{5}}\left[\left(\frac{1 + \sqrt{5}}{2}\right)^n - \left(\frac{1 - \sqrt{5}}{2}\right)^n\right]
        \]
        for each nonnegative integer $n$.
        \item Use (d) to conclude that if $n$ is a nonnegative integer, then the Fibonacci number $F_n$ is the integer that is closest to
        \[
            \frac{1}{\sqrt{5}}\left(\frac{1 + \sqrt{5}}{2}\right)^n.
        \]
    \end{subexercises}
    \textit{Each $F_n$ is a nonnegative integer, even though the right side of the formula in (d) does not look like an integer.
    The number}
    \[
        \frac{1 + \sqrt{5}}{2}
    \]
    \textit{is called the \emph{\textbf{golden ratio}}.}
\end{exercise}

\begin{exercise}[5d22]
    Suppose $T \in \lin{V}$ and $A$ is an $n$-by-$n$ matrix that is the matrix of $T$ with respect to some basis of $V$.
    Prove that if
    \[
        |A_{j, j}| > \sum_{\substack{k=1 \\ k \neq j}}^{n} |A_{j, k}|
    \]
    for each $j \in \set{1, \ldots, n}$, then $T$ is invertible.
    \textit{This exercise states that if the diagonal entries of the matrix of $T$ are large compared to the nondiagonal entries, then $T$ is invertible.}
\end{exercise}

\begin{exercise}[5d23]
    Suppose the definition of the Gershgorin disks is changed so that the radius of the $k^{\text{th}}$ disk is the sum of the absolute values of the entries in column (instead of row) $k$ of $A$, excluding the diagonal entry.
    Show that the Gershgorin disk theorem (5.67) still holds with this changed definition.
\end{exercise}